\subsection{The unit-connector eighth moment for every triple \((2,\allowbreak b,\allowbreak c)\), \(2<\allowbreak b<\allowbreak c\)}\label{proof:Q5-C070}

\paragraph{Q70.0 Statement and directions.}
For every finite probability space and real skew kernel \(|K|\le\allowbreak 1\),
and all finite integers \(2<\allowbreak b<\allowbreak c\), consider
\begin{equation}\tag{Q70.0.1}\label{eq:Q70:0.1}
F_D=E_0\bar f_2h_bf_c,\qquad H_D=E_rpT(F_D),
\end{equation}
where \(E_0,\allowbreak E_r\) are arbitrary finite genuine equal-pair products.
Then the directions specified below satisfy
\begin{equation}\tag{Q70.0.2}\label{eq:Q70:0.2}
\boxed{\E H_D^8\le e^{-d/27}\le1-d/28.}
\end{equation}
The entire two-edge arm points toward the terminal, giving
\(\bar f_2\). The \(b\)-arm has its first edge inward and its
remaining edges outward, giving \(h_b\); the \(c\)-arm is outward,
giving \(f_c\). The unit connector points from terminal to root,
whose unit arm points outward. Each equal pair uses an outward arm
and an arm first inward and then outward. Any finite number of pairs
and repeated lengths is allowed, with empty product one.
For \(b=\allowbreak 3,\allowbreak 4\), this direction differs from the corresponding
positive-step weighted-fourth direction of Q4-C143.

\paragraph{Q70.1 Finite model and short paths.}
Use the finite-model definitions in Appendix~\ref{proof:Q5-C064}:
\begin{equation}\tag{Q70.1.1}\label{eq:Q70:1.1}
\begin{gathered}
Jv=(\E v)\one,\quad A=J+K,\quad T=J-K=A^*,\quad
g=K\one,\quad p=1+g,\quad q=1-g,\\
f_n=A^n\one,\quad\bar f_n=T^n\one,\quad h_n=Tf_{n-1},
\quad r=Kg,\quad s=Kr,\\
\delta_j=\|K^j\one\|_2^2,\quad d=\delta_1,\quad
\theta=13/32,\quad\kappa=19/32,\quad\tau=d-\delta_2.
\end{gathered}
\end{equation}
After removing zero-mass atoms, all expectations, inner products,
spectra and adjoints use the same measure. The kernel is real,
skew, and bounded in absolute value by one. Its bounds are
\begin{equation}\tag{Q70.1.2}\label{eq:Q70:1.2}
0\le d\le1/3,\quad\|K\|^2\le\theta,\quad
|Kv(x)|\le\|v\|_2,\quad d^2\le\delta_2\le\theta d,\quad
\delta_{j+1}\le\theta\delta_j.
\end{equation}
For \(-K^2\), the spectral measure \(\nu\) of \(\one\) is
supported on \([0,\allowbreak \theta]\), with \(\delta_j=\allowbreak \int x^j\,\allowbreak d\nu(x)\).
Both \(A,\allowbreak T\) preserve order and \(p(x),\allowbreak q(x)\ge\allowbreak \mu(x)>\allowbreak 0\).
The exact short identities are
\begin{equation}\tag{Q70.1.3}\label{eq:Q70:1.3}
\E r=-d,\quad\E s=0,\quad\E(gr)=0,\quad
\bar f_2=q+r,\qquad\bar f_3=q+r-s-d.
\end{equation}
The arm and equal-pair bounds are
\begin{equation}\tag{Q70.1.4}\label{eq:Q70:1.4}
0\le f_n,\bar f_n,h_n\le2\quad(n\ge1),\qquad
0\le f_\ell h_\ell\le1.
\end{equation}
Thus \(0\le\allowbreak  E_0,\allowbreak E_r\le\allowbreak 1\).

\paragraph{Q70.2 A finite exponential bound for the long-arm pair.}
Put
\[
S_n=\sum_{j=0}^nK^j\one,\qquad
G_m(z)=\sum_{j=0}^{m-1}z^j,\qquad G_0=0,
\]
and
\(\beta_1=\allowbreak \beta_2=\allowbreak 0\),
\(\beta_n=\allowbreak \sum_{j=\allowbreak 1}^{\lfloor(n-1)/2\rfloor}(-1)^{j-1}\delta_j\).
The finite convolution and geometric bounds are
\begin{equation}\tag{Q70.2.1}\label{eq:Q70:2.1}
S_n=f_n+\sum_{j=3}^n\beta_jf_{n-j},\qquad
\tau\le\beta_n\le d\quad(n\ge3),
\end{equation}
\begin{equation}\tag{Q70.2.2}\label{eq:Q70:2.2}
\|G_m(K)\|\le1\quad(m=0,1\text{ or }m\ge3),\qquad
\|G_2(K)\|^2\le1+\theta.
\end{equation}
For \(3\le\allowbreak  b<\allowbreak c\), set
\begin{equation}\tag{Q70.2.3}\label{eq:Q70:2.3}
R=R_{bc}=\sum_{j=b}^{c-1}K^jg=K^bG_{c-b}(K)g,\qquad
\Gamma=\begin{cases}2d-\delta_2,&b=3,\\3d,&b\ge4.\end{cases}
\end{equation}
The pair satisfies
\begin{equation}\tag{Q70.2.4}\label{eq:Q70:2.4}
0\le h_bf_c\le e^{R-\Gamma}.
\end{equation}
The case \(b=\allowbreak 3\) is the finite short-pair inequality of Q3-C014,
Lemma 6.1, also displayed as (Q78.1.7) in
Appendix~\ref{proof:Q5-C078}.
For \(b\ge\allowbreak 4\), subtract the finite convolutions:
\begin{equation}\tag{Q70.2.5}\label{eq:Q70:2.5}
f_c-f_b=R+
\sum_{i=0}^{b-3}(\beta_{b-i}-\beta_{c-i})f_i
-\sum_{i=b-2}^{c-3}\beta_{c-i}f_i.
\end{equation}
The positive part of the first sum is at most
\(4\delta_2/(1-\theta^2)\le\allowbreak 2d\).
For every \(n\ge\allowbreak 4\), the finite path-mean bound is
\(\E f_n\le\allowbreak 1-83d/32\le\allowbreak 1-5d/2\), proved in
Appendix~\ref{proof:Q5-C090}. Hence for \(b\ge\allowbreak 5\),
\(h_b+f_c=\allowbreak 2\E f_{b-1}+f_c-f_b\le\allowbreak 2+R-3d\).
For \(b=\allowbreak 4\), retain the \(i=\allowbreak 2\) term in the last sum, including
the shortest case \(c=\allowbreak 5\):
\fitdisplay{f_c-f_4\le R+\delta_2(1+p)-\tau f_2.}{}
The row bound \(|r|\le\allowbreak  pq/2\) gives \(f_2=\allowbreak p+r\ge\allowbreak  p^2/2\).
Completing the square for \(p\ge\allowbreak 0\), and using
\(\delta_2\le\allowbreak  d/2\), \(\tau\ge\allowbreak  d/2\), gives
\fitdisplay{
\delta_2(1+p)-\tau f_2
\le\delta_2+\delta_2p-\tau p^2/2
\le\delta_2+\frac{\delta_2^2}{2\tau}\le3d/4.
}{}
Thus \(h_4+f_c\le\allowbreak 2+R-13d/4\le\allowbreak 2+R-3d\).
For \(d>\allowbreak 0\), \(\tau>\allowbreak 0\); the zero boundary is treated below.
The arm sum is nonnegative, so AM--GM gives
\(h_bf_c\le\allowbreak [(h_b+f_c)/2]^2\le\allowbreak  e^{R-\Gamma}\).
Let \(J_R=\allowbreak \E R^2\). Equation~\eqref{eq:Q70:2.2} gives
\begin{equation}\tag{Q70.2.6}\label{eq:Q70:2.6}
|R|^2\le(1+\theta)\theta^{b-1}d,\qquad
J_R\le(1+\theta)\theta^bd.
\end{equation}
In particular,
\begin{equation}\tag{Q70.2.7}\label{eq:Q70:2.7}
|R|^2\le7605d/32768\le7605/98304<9/100,\qquad |R|<3/10,
\end{equation}
where \(760500<\allowbreak 884736\).

\paragraph{Q70.3 Propagation retaining the two-edge arm.}
The negative-mean lemma \eqref{eq:Q65:2.4} applies whenever
\(\E N\le\allowbreak 0\), \(|N|\le\allowbreak \sqrt{A_0d}\),
\(\E N^2\le\allowbreak \theta A_0d\), and \(|\E N|\le\allowbreak \sqrt{A_0}d\):
\begin{equation}\tag{Q70.3.1}\label{eq:Q70:3.1}
\E e^{8N}\le\exp\{8\E N+\lambda(\sqrt d)A_0d\},\qquad
\lambda(t)=\frac{188}{25}\left(2-\frac{\kappa}{1+t}\right)^2.
\end{equation}
Since \(\lambda\) increases and \(\sqrt d\le\allowbreak 1/\sqrt3<\allowbreak 3/5\),
\begin{equation}\tag{Q70.3.2}\label{eq:Q70:3.2}
\lambda(\sqrt d)\le\lambda(3/5)=32691132/1638400<20.
\end{equation}
Here \(2-\kappa/(8/5)=\allowbreak 417/256\), and
\(188\cdot417^2=\allowbreak 32691132<\allowbreak 32768000=\allowbreak 20\cdot1638400\).
For every real \(u\le\allowbreak 3/10\),
\begin{equation}\tag{Q70.3.3}\label{eq:Q70:3.3}
e^u\le1+u+(3/5)u^2.
\end{equation}
Indeed, \((e^u-1-u)/u^2=\allowbreak \int_0^1(1-t)e^{tu}\,\allowbreak dt\) increases
with \(u\). At \(u=\allowbreak 3/10\), the exponential series and
\(j!\ge\allowbreak 2\cdot3^{j-2}\) give
\fitdisplay{
e^u\le1+u+\frac{u^2}{2(1-u/3)}=27/20<677/500
=1+3/10+(3/5)(3/10)^2.
}{}
Use the continuous extension at zero. Set
\begin{equation}\tag{Q70.3.4}\label{eq:Q70:3.4}
\begin{gathered}
G=g-r,\quad\bar f_2=1-G,\quad\E G^2=d+\delta_2,\\
V=g-r-R,\quad N=KV=r-s-KR,\\
C=-d+\E R-\E(GR)+\sqrt{(d+\delta_2)J_R}+(12/5)J_R.
\end{gathered}
\end{equation}
Positivity of \(T\), \(\bar f_2\ge\allowbreak 0\), and \eqref{eq:Q70:3.3} give
\fitdisplay{
T(\bar f_2e^R)\le\bar f_3+T(\bar f_2R)+(3/5)T(\bar f_2R^2).
}{}
The exact linear term is
\(T(\bar f_2R)=\allowbreak \E R-\E(GR)-KR+K(GR)\), and
\fitdisplay{
K(GR)\le\E|GR|\le\sqrt{(d+\delta_2)J_R},\qquad
T(\bar f_2R^2)\le2\E(\bar f_2R^2)\le4J_R.
}{}
Therefore
\begin{equation}\tag{Q70.3.5}\label{eq:Q70:3.5}
T(\bar f_2e^R)\le q+C+N.
\end{equation}
The right side is nonnegative. By \(p+q=\allowbreak 2\) and AM--GM,
\(p(q+C+N)\le\allowbreak [1+(C+N)/2]^2\le\allowbreak  e^{C+N}\), with nonnegative base.
For the undecorated \(H=\allowbreak pT(\bar f_2h_bf_c)\),
\begin{equation}\tag{Q70.3.6}\label{eq:Q70:3.6}
0\le H\le e^{-\Gamma+C+N}.
\end{equation}
The coupled mean identities are
\begin{equation}\tag{Q70.3.7}\label{eq:Q70:3.7}
\begin{gathered}
\E N=-d+\E(gR),\\
C+\E N=-2d+\E R+\E(rR)
                +\sqrt{(d+\delta_2)J_R}+(12/5)J_R.
\end{gathered}
\end{equation}
Thus \(\E(gR)\) cancels between the terms.
If \(\|V\|^2\le\allowbreak  A_0d\) and \(\E N\le\allowbreak 0\), the row and spectral
bounds and Cauchy give all hypotheses of \eqref{eq:Q70:3.1}.
Using \eqref{eq:Q70:3.2} then gives
\begin{equation}\tag{Q70.3.8}\label{eq:Q70:3.8}
\boxed{\begin{aligned}
\log\E H^8\le{}&-8\Gamma-16d+8\{\E R+\E(rR)\}\\
 &+8\sqrt{(d+\delta_2)J_R}+(96/5)J_R+20A_0d.
\end{aligned}}
\end{equation}

\paragraph{Q70.4 Mixed moments and the negative mean.}
Put \(a_j=\allowbreak \E K^j\one\), so \(a_{2j}=\allowbreak (-1)^j\delta_j\) and
odd \(a_j\) vanish. Skewness gives
\begin{equation}\tag{Q70.4.1}\label{eq:Q70:4.1}
\E(rK^jg)=a_{j+3},\qquad\E(K^jg)=a_{j+1}.
\end{equation}
Write \(e_n=\allowbreak \delta_n-\delta_{n+1}\). The spectral representation gives
\begin{equation}\tag{Q70.4.2}\label{eq:Q70:4.2}
e_n=\int x^n(1-x)\,d\nu(x)\ge0,\qquad e_{n+1}\le\theta e_n.
\end{equation}
For each odd \(j=\allowbreak 2n-1\), the sum in \eqref{eq:Q70:4.1} is
\((-1)^ne_n\); for even \(j\), it is zero.
The contiguous finite interval in \(R\) therefore gives a finite
alternating sum of decreasing nonnegative terms, and
\begin{equation}\tag{Q70.4.3}\label{eq:Q70:4.3}
\E R+\E(rR)\le
\begin{cases}\delta_2-\delta_3,&b=3,\\0,&b=4,\\\delta_4,&b\ge5.\end{cases}
\end{equation}
The first class begins with \(e_2\), and equals
\(\delta_2-\delta_3\) for \((b,\allowbreak c)=\allowbreak (3,\allowbreak 4),\allowbreak (3,\allowbreak 5)\).
The second, when nonempty, begins with \(-e_3\).
In the third, a positive first term has index at least four and is
at most \(e_4\le\allowbreak \delta_4\); a negative first term gives upper bound zero.
Similarly, only even \(j\) contribute to \(\E(gK^jg)\). For \(j\ge\allowbreak 3\),
the first possible positive term is \(j=\allowbreak 4\), with value \(\delta_3\),
followed by decreasing alternating terms. Every finite truncation gives
\begin{equation}\tag{Q70.4.4}\label{eq:Q70:4.4}
\E(gR)\le\delta_3\le\theta^2d<d\qquad(d>0).
\end{equation}
Hence \(\E N<\allowbreak 0\) in every class.

\paragraph{Q70.5 Norms for every finite \(b,\allowbreak c\).}
On a rotation block write \(z=\allowbreak i\sqrt x\), \(0\le\allowbreak  x\le\allowbreak \theta\).
Then \(V=\allowbreak P(K)g\), with
\begin{equation}\tag{Q70.5.1}\label{eq:Q70:5.1}
P(z)=1-z-\sum_{j=b}^{c-1}z^j.
\end{equation}
For \((b,\allowbreak c)=\allowbreak (3,\allowbreak 4)\), \(P=\allowbreak 1-z-z^3\), so
\begin{equation}\tag{Q70.5.2}\label{eq:Q70:5.2}
|P(i\sqrt x)|^2=1+x(1-x)^2\le1+4/27=31/27.
\end{equation}
On \([0,\allowbreak 1]\), the derivative of \(x(1-x)^2\) is
\((1-x)(1-3x)\), giving maximum \(4/27\). Here \(J_R=\allowbreak \delta_4\).
For \((b,\allowbreak c)=\allowbreak (3,\allowbreak 5)\), \(P=\allowbreak 1-z-z^3-z^4\), and
\fitdisplay{|P|^2=(1-x^2)^2+x(1-x)^2=1+x-4x^2+x^3+x^4.}{}
Since \(x+x^2\le\allowbreak \theta+\theta^2<\allowbreak 1\),
\begin{equation}\tag{Q70.5.3}\label{eq:Q70:5.3}
|P|^2\le1+x-3x^2\le13/12.
\end{equation}
This case has the exact second moment \(J_R=\allowbreak \delta_4+\delta_5\),
retaining the exceptional \(G_2\) contribution.
For \(b=\allowbreak 3,\allowbreak c\ge\allowbreak 6\), the finite sums satisfy
\fitdisplay{
x^2-x^3\le\operatorname{Re}\sum_{j=3}^{c-1}z^j\le x^2,
\qquad
-x^{3/2}\le\operatorname{Im}\sum_{j=3}^{c-1}z^j
\le-x^{3/2}+x^{5/2}.
}{}
Thus \(\operatorname{Re}P>\allowbreak 0\), \(\operatorname{Im}P\le\allowbreak 0\), and
\begin{equation}\tag{Q70.5.4}\label{eq:Q70:5.4}
\begin{aligned}
|P|^2&\le(1-x^2+x^3)^2+x(1-x+x^2)^2\\
 &=1+x-4x^2+5x^3-x^4-x^5+x^6\\
 &\le1+x-(63/32)x^2\le1+8/63=71/63.
\end{aligned}
\end{equation}
Use \(-x^4-x^5+x^6\le\allowbreak 0\), \(4-5\theta=\allowbreak 63/32\).
Here \(c-3\ge\allowbreak 3\), so \eqref{eq:Q70:2.2} gives \(J_R\le\allowbreak \delta_4\).
For \(b=\allowbreak 4,\allowbreak c>\allowbreak 4\), the contiguous sum has real part in
\([x^2-x^3,\allowbreak x^2]\) and imaginary part in \([0,\allowbreak x^{5/2}]\). Therefore
\fitdisplay{
|P|^2\le(1-x^2+x^3)^2+x(1+x^2)^2
=1+x-2x^2+4x^3+x^4-x^5+x^6.
}{}
As \(x\le\allowbreak 1/2\), both \(-2x^2+4x^3\) and \(-x^5+x^6\) are nonpositive,
so
\begin{equation}\tag{Q70.5.5}\label{eq:Q70:5.5}
|P|^2\le1+x+x^4<1+13/32+1/32=23/16<3/2.
\end{equation}
Here \(\theta^4=\allowbreak 28561/1048576<\allowbreak 1/32\), and
\eqref{eq:Q70:2.6} also gives \(J_R<\allowbreak (3/64)d\).
For \(b\ge\allowbreak 5,\allowbreak c>\allowbreak b\), the finite alternating sums give
\fitdisplay{
-x^3\le\operatorname{Re}\sum_{j=b}^{c-1}z^j\le x^4,\qquad
\left|\operatorname{Im}\sum_{j=b}^{c-1}z^j\right|\le x^{5/2}.
}{}
Since \(\operatorname{Re}P>\allowbreak 0\),
\begin{equation}\tag{Q70.5.6}\label{eq:Q70:5.6}
|P|^2\le(1+x^3)^2+x(1+x^2)^2
=1+x+4x^3+x^5+x^6<7/4.
\end{equation}
For the rational comparison use \(x\le\allowbreak \theta<\allowbreak 5/12\),
\(4\theta^3<\allowbreak 3/10\), and
\(\theta^5+\theta^6<\allowbreak (3/2)/90<\allowbreak 1/30\); these sum to less than \(3/4\).
The cross-products for \(\theta^5<\allowbreak 1/90\) and \(4\theta^3<\allowbreak 3/10\)
are \(33416370<\allowbreak 33554432\) and \(87880<\allowbreak 98304\).
Also
\begin{equation}\tag{Q70.5.7}\label{eq:Q70:5.7}
J_R\le(1+\theta)\theta^5d<d/60.
\end{equation}
Integrating the block bounds gives \(\E V^2\le\allowbreak  A_0d\) in each class.

\paragraph{Q70.6 The eighth-moment bounds on \(0<\allowbreak d\le\allowbreak 1/3\).}
Apply \eqref{eq:Q70:3.8}, retaining the common contribution
\(\E R+\E(rR)+(12/5)J_R\).
For \((b,\allowbreak c)=\allowbreak (3,\allowbreak 4)\), \(J_R=\allowbreak \delta_4\le\allowbreak \theta\delta_3\), and
\begin{equation}\tag{Q70.6.1}\label{eq:Q70:6.1}
\E R+\E(rR)+(12/5)J_R
\le\delta_2-\delta_3+(12/5)\delta_4
\le\delta_2-\delta_3/40\le\delta_2.
\end{equation}
Moreover,
\begin{equation}\tag{Q70.6.2}\label{eq:Q70:6.2}
\sqrt{(d+\delta_2)J_R}
\le\sqrt{(1+\theta)\theta^3}\,d<5d/16.
\end{equation}
The squared comparison is \(98865<\allowbreak 102400\), with denominator
\(1048576\). Use \(\Gamma=\allowbreak 2d-\delta_2\), \(A_0=\allowbreak 31/27\):
\begin{equation}\tag{Q70.6.3}\label{eq:Q70:6.3}
\log\E H^8\le(-32+16\theta+5/2+20\cdot31/27)d=-d/27.
\end{equation}
The final equality is \(-23+620/27=\allowbreak -1/27\).
For \((b,\allowbreak c)=\allowbreak (3,\allowbreak 5)\), use the exact \(J_R=\allowbreak \delta_4+\delta_5\):
\begin{equation}\tag{Q70.6.4}\label{eq:Q70:6.4}
\begin{aligned}
\E R+\E(rR)+(12/5)J_R
&\le\delta_2+[-1+(12/5)\theta(1+\theta)]\delta_3\\
&=\delta_2+(95/256)\delta_3<\delta_2+d/16.
\end{aligned}
\end{equation}
The last comparison is \(16055<\allowbreak 16384\), with denominator \(262144\).
Also
\begin{equation}\tag{Q70.6.5}\label{eq:Q70:6.5}
\sqrt{(d+\delta_2)J_R}\le(1+\theta)\theta^{3/2}d<3d/8,
\end{equation}
where the squared comparison is \(4448925<\allowbreak 4718592\), with
denominator \(33554432\). Using \(A_0=\allowbreak 13/12\),
\begin{equation}\tag{Q70.6.6}\label{eq:Q70:6.6}
\log\E H^8\le(-32+16\theta+1/2+3+20\cdot13/12)d=-d/3.
\end{equation}
For \(b=\allowbreak 3,\allowbreak c\ge\allowbreak 6\), \(J_R\le\allowbreak \delta_4\), so
\eqref{eq:Q70:6.1}--\eqref{eq:Q70:6.2} apply. With \(A_0=\allowbreak 71/63\),
\begin{equation}\tag{Q70.6.7}\label{eq:Q70:6.7}
\log\E H^8\le(-23+20\cdot71/63)d=-(29/63)d.
\end{equation}
For \(b=\allowbreak 4,\allowbreak c>\allowbreak 4\), \(\Gamma=\allowbreak 3d\), \(\E R+\E(rR)\le\allowbreak 0\),
\(A_0=\allowbreak 3/2\). From \(J_R<\allowbreak 3d/64\), \(d+\delta_2<\allowbreak (3/2)d\),
\fitdisplay{\sqrt{(d+\delta_2)J_R}<d/3,\qquad (12/5)J_R<d/8.}{}
The comparisons are \(9/128<\allowbreak 1/9\) and \(9/80<\allowbreak 1/8\). Hence
\begin{equation}\tag{Q70.6.8}\label{eq:Q70:6.8}
\log\E H^8\le(-40+30+8/3+1)d=-(19/3)d.
\end{equation}
For \(b\ge\allowbreak 5,\allowbreak c>\allowbreak b\), \(\Gamma=\allowbreak 3d\), \(A_0=\allowbreak 7/4\), and
\(\E R+\E(rR)\le\allowbreak \delta_4<\allowbreak d/14\), using \(\theta^3<\allowbreak 1/14\)
and \(30758<\allowbreak 32768\). Since \(J_R<\allowbreak d/60\),
\fitdisplay{\sqrt{(d+\delta_2)J_R}<d/6,\qquad (12/5)J_R<d/25.}{}
Now \(1/14+1/6+1/25=\allowbreak 146/525<\allowbreak 3/10\), giving
\begin{equation}\tag{Q70.6.9}\label{eq:Q70:6.9}
\log\E H^8\le(-40+35+12/5)d=-13d/5.
\end{equation}
The five classes cover every \(3\le\allowbreak  b<\allowbreak c\). Their common bound is
\(\E H^8\le\allowbreak  e^{-d/27}\).

\paragraph{Q70.7 Boundary, decorations and full reversal.}
At \(d=\allowbreak 0\), \(g=\allowbreak 0\), \(A\one=\allowbreak T\one=\allowbreak \one\), and leaf integration
makes every rooted tree message one, although \(K\) may be nonzero.
For the decorations, positivity gives
\fitdisplay{0\le E_rpT(E_0\bar f_2h_bf_c)\le pT(\bar f_2h_bf_c).}{}
Full reversal replaces \(K\) by \(-K\), exchanges \(A/T\), \(p/q\),
and changes \(\bar f_2\) to \(f_2\), \(h_b\) to
\(W_b=\allowbreak A\bar f_{b-1}\), and \(f_c\) to \(\bar f_c\), with every
decoration reversed as well. The value of \(d\) is unchanged.
Finally,
\fitdisplay{e^{-d/27}\le\frac1{1+d/27}
=1-\frac d{27+d}\le1-d/28\qquad(d\le1/3).}{}
This proves \eqref{eq:Q70:0.2}.
For the direction comparison, Q4-C143, \S5.3, uses
\(\bar f_2f_3h_c\) at \(b=\allowbreak 3\) and \(\bar f_2f_4h_c\) at \(b=\allowbreak 4\),
whereas the present input is \(\bar f_2h_bf_c\).
For \(b\ge\allowbreak 5\), the directions agree, so the existing all-positive-length
\(p\)-weighted fourth and the present unit eighth can be used for the
same input.
